\section{Introduction}\label{sec:introduction}

Conventional therapeutic modalities remain largely ineffective against many disease-relevant proteins, particularly those mediating protein-protein interactions (PPIs). Small molecules are inherently restricted to targeting deep, well-defined binding pockets and cannot effectively engage the large, shallow, and dynamic interfaces characteristic of PPIs \citep{bojadzic2018toward}. Monoclonal antibodies, while capable of high-affinity binding to these extended surfaces, suffer from obligate parenteral administration and are unable to access intracellular targets.
Macrocyclic peptides occupy a unique therapeutic niche. Their conformationally constrained backbones present preorganized interaction surfaces that match or exceed the footprint of antibody complementarity-determining regions (CDRs), enabling antibody-like potency and selectivity \citep{zorzi2017cyclic}. Critically, rational chemical optimization can confer clinically meaningful oral bioavailability, as demonstrated by the late-stage oral PCSK9 inhibitor enlicitide (MK-0616) \citep{gare2025lead}.


Despite remarkable advances in generating structures of target-binding cyclic peptides \citep{rettie2025accurate,rettie2025cyclic}, existing structural models are predominantly trained on natural amino acid proteins and peptides, limiting their ability to design cyclic peptides incorporating non-natural amino acids. Non-natural amino acids are indispensable for enhancing proteolytic stability, augmenting binding affinity via novel intermolecular interactions, and providing versatile attachment points for intramolecular cyclization \citep{hickey2023beyond}. Therefore, the development of generalizable foundation models for cyclic peptides that can accommodate non-natural amino acids and diverse cyclization chemistries is of paramount importance.

% Many disease-relevant proteins remain difficult to modulate with conventional small molecules or biologics because their functional surfaces are large, shallow, dynamic, or intracellular. Macrocyclic peptides provide an attractive therapeutic middle ground: their constrained backbones can present extended interaction surfaces with antibody-like specificity while remaining small enough to exploit passive or assisted cellular entry in selected chemical regimes \citep{ref_buyanova_pei}. The same features that make them compelling also make them hard to design. Productive macrocycles often combine natural amino acids, non-natural amino acids, stereochemical variants, backbone modifications, and intramolecular cyclization bonds; a useful computational generator must therefore reason jointly over sequence, residue chemistry, synthetic attachment points, and pharmacological objectives such as permeability and target binding.

The simplified molecular-input line-entry system (SMILES) remains a dominant representation for molecular machine learning because it encodes atoms and bonds in a compact one-dimensional string \citep{ref_smiles}. Recent peptide generators have therefore explored SMILES-space modeling. Geylan et al. present PepINVENT \citep{geylan2025pepinvent}, a transformer-based generative AI framework for SMILES generation, which enables de novo peptide design with both natural and non-natural amino acids and supports diverse topologies ranging from linear peptides to multiple macrocyclic architectures. Trained on a text-infilling task and coupled with reinforcement learning, PepINVENT facilitates multi-parameter optimization of peptides for therapeutic properties such as solubility, permeability, and topological constraints. Tang et al. developed PepTune \citep{ref_peptune}, which formulates therapeutic peptide design as a multi-objective guided discrete diffusion process over peptide SMILES. While SMILES can precisely represent complex macrocyclic molecules and broadly compatible with existing molecular property predictors, it encounters significant scalability challenges when applied to macrocyclic peptides. A single residue-level design decision may correspond to many coordinated SMILES tokens describing backbone atoms, side-chain atoms, stereochemistry, branches, and ring closures. Because available macrocyclic peptide datasets are far smaller than the resulting character-level state space, atom-by-atom generation can make validity learning data-hungry and can complicate reinforcement learning or guided diffusion, where sparse property rewards must be propagated across long token sequences.

Residue-level notation is therefore a natural bridge from molecular strings to peptide design, and HELM provides a standard abstraction. The hierarchical editing language for macromolecules is widely used for peptide and biopolymer therapeutics, including macrocyclic peptides, in databases and registration workflows \citep{ref_helm,ref_chembl}. HELM defines reusable monomers, polymer order, and inter-monomer bonds through explicit attachment points, encoding natural residues, non-natural residues, modifications, and cyclizations as monomer symbols plus typed R-group connections. This modular grammar suits macromolecules by retaining sequence compactness and covalent topology, while matching residue-level resources such as ChEMBL and CycPeptMPDB for data-driven macrocyclic peptide modeling \citep{ref_chembl,ref_cycpeptmpdb}.

HELM-GPT, a GPT-style decoder, has established HELM as a learnable language for \textit{de novo} macrocyclic peptide design and property optimization \citep{ref_helm_gpt}. However, a purely string-level HELM language model still treats monomer symbols largely as textual tokens, making its representations dependent on the limited peptide-sequence corpora available in this domain rather than on the underlying molecular structures of the monomers. Consequently, the explored chemical space is not guaranteed to extend beyond neighborhoods of the corpus distribution, which can leave rare building blocks under-explored, even when their physicochemical profiles closely resemble those of common monomers.

In this study, we propose \mytool, an ALD framework for \textit{de novo} macrocyclic peptide generation. PepALD adopts the HELM tokenization strategy, but embeds each monomer with structured representations derived from the Uni-Mol chemical foundation model \citep{ref_unimol}, so that generation is performed in a chemically informed latent space rather than over symbolic tokens alone. At each autoregressive step, a causal context encoder summarizes the generated prefix, a context-conditioned diffusion module samples the next monomer embedding, and a constrained mapper projects it back to an admissible HELM building block \citep{ref_ddpm,ref_ddim}. In parallel, an R-group-aware ring bond predictor reuses the autoregressive context together with pretrained attachment-site embeddings to determine whether, where, and how the peptide should cyclize. To further steer PepALD toward desired properties, we build on diffusion-adapted direct preference optimization \citep{ref_dpo,ref_diffusion_dpo} and introduce a winner-protected Diffusion-DPO (WP-DPO) objective inspired by DPO-Positive \citep{ref_dpop}, while constructing preference pairs with a structure-aware strategy that combines maximal marginal relevance diversity selection with nearest hard-negative coupling \citep{ref_mmr}.

We conduct extensive experiments comparing PepALD with representative baselines, including HELM-GPT and PepTune \citep{ref_helm_gpt,ref_peptune}. The model is first pre-trained on ChEMBL32 HELM peptide sequences to learn general peptide syntax and monomer usage patterns, and then fine-tuned on CycPeptMPDB to capture macrocyclic peptide chemistry and topology \citep{ref_chembl,ref_cycpeptmpdb}. The resulting ALD model exhibits state-of-the-art generation performance, with leading validity, uniqueness, diversity, and novelty, as well as a balanced nearest-neighbor similarity profile. As target-focused case studies, we design cyclic peptide inhibitors against two protein targets: the intracellular SPSB2--iNOS protein--protein interaction \citep{sadek2018cyclic} and Mycobacterium tuberculosis chorismate mutase (MtbCM) \citep{van2025active}, using a cyclic peptide permeability predictor and Uni-Dock/Vina docking scores as reward components \citep{ref_unidock}. Under the WP-DPO framework, the weighted reward improves both permeability-oriented and affinity-oriented objectives. Benefiting from its chemically grounded latent space, ring-closure predictor, and preference-pair construction, PepALD demonstrates broad application potential in \textit{de novo} cyclic peptide design scenarios.
